% GENERATED FILE. Edit canonical lessons, then run npm run generate:books.
% canonical-source-hash: fnv1a64:3ae76f36171d71de
% canonical-lessons: ML-C369-vidhi, ML-C369-kurrakkaran, ML-C369-jamyam, ML-C369-niraparadhi, ML-C369-vidhikkuka

\chapter{The Verdict}
\label{ch:ml-the-verdict}
\hlchaptermodality{369}

\begin{chapteropening}
\textbf{By the end of this chapter:} \emph{I can say a verdict, the guilty one, bail, an innocent person and to pass sentence.}

Complete the last lesson of chapter 369: I can say a verdict, the guilty one, bail, an innocent person and to pass sentence.
\end{chapteropening}

\section[\texorpdfstring{\ml{വിധി} (vidhi)}{വിധി (vidhi)}]{\ml{വിധി} (vidhi) — a verdict; fate}

\label{lesson:ML-C369-vidhi}



\begin{quote}
Before the new one: say the Malayalam for a judge, then the Malayalam for the accused.
\end{quote}

\subsection*{You'll want to know: \ml{വിധി}}
\textbf{\ml{വിധി}} — \emph{vidhi} — \textquotedblleft{}a verdict; fate\textquotedblright{}.

\begin{quote}
\ml{കോടതിയുടെ} \ml{വിധി} \ml{ഇന്ന്} \ml{വരും}.
\end{quote}

\begin{quote}
\emph{kōṭatiyuṭe vidhi innŭ varuṁ.}
\end{quote}

\begin{quote}
The court's verdict will come today.
\end{quote}

Also fate, destiny: \ml{അത്} \ml{എന്റെ} \ml{വിധി} \ml{ആയിരുന്നു}, that was my fate.

\begin{cousinweb}
\textbf{\ml{വിധി}} (\emph{vidhi}) is borrowed from Sanskrit \emph{vidhi}, \textquotedblleft{}rule\textquotedblright{}, and later \textquotedblleft{}fate\textquotedblright{} (CDIAL 11755). A verdict is the rule a court lays down for one case.
\end{cousinweb}

\begin{grammarlens}[title={What you've built}]
One more word for telling what happened.
\end{grammarlens}

\subsection*{Your turn}
\begin{itemize}
\raggedright
  \item \emph{Say it:} \emph{vidhi}
  \item \emph{Say it:} \emph{vidhi}, once more
  \item \emph{Say it:} \emph{kōṭatiyuṭe vidhi innŭ varuṁ.} — The court's verdict will come today.
  \item \emph{Recall:} say the Malayalam for a judge, then the Malayalam for the accused, then say \emph{vidhi} again
\end{itemize}

\begin{tcolorbox}[breakable,colback=teal!4,colframe=teal!35!black,title={Before you move on}]
What does \ml{വിധി} mean? (\textquotedblleft{}A verdict; fate\textquotedblright{}.) Say it once more.
\end{tcolorbox}
\section[\texorpdfstring{\ml{കുറ്റക്കാരൻ} (kuṟṟakkāran)}{കുറ്റക്കാരൻ (kuṟṟakkāran)}]{\ml{കുറ്റക്കാരൻ} (kuṟṟakkāran) — the guilty one; a culprit}

\label{lesson:ML-C369-kurrakkaran}



\begin{quote}
Before the new one: say the Malayalam for the accused, then the Malayalam for a verdict.
\end{quote}

\subsection*{You'll want to know: \ml{കുറ്റക്കാരൻ}}
\textbf{\ml{കുറ്റക്കാരൻ}} — \emph{kuṟṟakkāran} — \textquotedblleft{}the guilty one; a culprit\textquotedblright{}.

\begin{quote}
\ml{കോടതി} \ml{അവനെ} \ml{കുറ്റക്കാരനായി} \ml{കണ്ടെത്തി}.
\end{quote}

\begin{quote}
\emph{kōṭati avane kuṟṟakkāranāyi kaṇṭetti.}
\end{quote}

\begin{quote}
The court found him guilty.
\end{quote}

Not guilty is \ml{കുറ്റക്കാരൻ} \ml{അല്ല}. -\ml{ക്കാരൻ} makes a person from a noun.

\begin{cousinweb}
\textbf{\ml{കുറ്റക്കാരൻ}} (\emph{kuṟṟakkāran}) is native Dravidian: \ml{കുറ്റം} (\emph{kuṟṟaṁ}), \textquotedblleft{}an offence\textquotedblright{} (DEDR 1851), with the ending -\ml{ക്കാരൻ}, \textquotedblleft{}one who has it or does it\textquotedblright{}.
\end{cousinweb}

\begin{grammarlens}[title={What you've built}]
One more word for telling what happened.
\end{grammarlens}

\subsection*{Your turn}
\begin{itemize}
\raggedright
  \item \emph{Say it:} \emph{kuṟṟakkāran}
  \item \emph{Say it:} \emph{kuṟṟakkāran}, once more
  \item \emph{Say it:} \emph{kōṭati avane kuṟṟakkāranāyi kaṇṭetti.} — The court found him guilty.
  \item \emph{Recall:} say the Malayalam for the accused, then the Malayalam for a verdict, then say \emph{kuṟṟakkāran} again
\end{itemize}

\begin{tcolorbox}[breakable,colback=teal!4,colframe=teal!35!black,title={Before you move on}]
What does \ml{കുറ്റക്കാരൻ} mean? (\textquotedblleft{}The guilty one; a culprit\textquotedblright{}.) Say it once more.
\end{tcolorbox}
\section[\texorpdfstring{\ml{ജാമ്യം} (jāmyaṁ)}{ജാമ്യം (jāmyaṁ)}]{\ml{ജാമ്യം} (jāmyaṁ) — bail}

\label{lesson:ML-C369-jamyam}



\begin{quote}
Before the new one: say the Malayalam for a verdict, then the Malayalam for the guilty one.
\end{quote}

\subsection*{You'll want to know: \ml{ജാമ്യം}}
\textbf{\ml{ജാമ്യം}} — \emph{jāmyaṁ} — \textquotedblleft{}bail\textquotedblright{}.

\begin{quote}
\ml{കോടതി} \ml{അവന്} \ml{ജാമ്യം} \ml{നൽകി}.
\end{quote}

\begin{quote}
\emph{kōṭati avanŭ jāmyaṁ nalki.}
\end{quote}

\begin{quote}
The court granted him bail.
\end{quote}

Bail is granted: \ml{ജാമ്യം} \ml{നൽകുക}. Out on bail is \ml{ജാമ്യത്തിൽ}.

\begin{cousinweb}
The origin of \textbf{\ml{ജാമ്യം}} (\emph{jāmyaṁ}) is not settled. It means a guarantee given to a court, so that the accused may go home until the trial.
\end{cousinweb}

\begin{grammarlens}[title={What you've built}]
One more word for telling what happened.
\end{grammarlens}

\subsection*{Your turn}
\begin{itemize}
\raggedright
  \item \emph{Say it:} \emph{jāmyaṁ}
  \item \emph{Say it:} \emph{jāmyaṁ}, once more
  \item \emph{Say it:} \emph{kōṭati avanŭ jāmyaṁ nalki.} — The court granted him bail.
  \item \emph{Recall:} say the Malayalam for a verdict, then the Malayalam for the guilty one, then say \emph{jāmyaṁ} again
\end{itemize}

\begin{tcolorbox}[breakable,colback=teal!4,colframe=teal!35!black,title={Before you move on}]
What does \ml{ജാമ്യം} mean? (\textquotedblleft{}Bail\textquotedblright{}.) Say it once more.
\end{tcolorbox}
\section[\texorpdfstring{\ml{നിരപരാധി} (niraparādhi)}{നിരപരാധി (niraparādhi)}]{\ml{നിരപരാധി} (niraparādhi) — an innocent person}

\label{lesson:ML-C369-niraparadhi}



\begin{quote}
Before the new one: say the Malayalam for the guilty one, then the Malayalam for bail.
\end{quote}

\subsection*{You'll want to know: \ml{നിരപരാധി}}
\textbf{\ml{നിരപരാധി}} — \emph{niraparādhi} — \textquotedblleft{}an innocent person\textquotedblright{}.

\begin{quote}
\ml{ഞാൻ} \ml{നിരപരാധിയാണ്} \ml{എന്ന്} \ml{അവൻ} \ml{പറഞ്ഞു}.
\end{quote}

\begin{quote}
\emph{ñān niraparādhiyāṇŭ ennŭ avan paṟaññu.}
\end{quote}

\begin{quote}
He said that he was innocent.
\end{quote}

\ml{നിർ}-, without, + \ml{അപരാധം}, an offence: one who did nothing wrong.

\begin{cousinweb}
\textbf{\ml{നിരപരാധി}} (\emph{niraparādhi}) is borrowed from Sanskrit \emph{niraparādha}, \textquotedblleft{}unoffending, guiltless\textquotedblright{} (Monier-Williams), made of \emph{nir-}, \textquotedblleft{}without\textquotedblright{}, and \emph{aparādha}, \textquotedblleft{}an offence\textquotedblright{}.
\end{cousinweb}

\begin{grammarlens}[title={What you've built}]
One more word for telling what happened.
\end{grammarlens}

\subsection*{Your turn}
\begin{itemize}
\raggedright
  \item \emph{Say it:} \emph{niraparādhi}
  \item \emph{Say it:} \emph{niraparādhi}, once more
  \item \emph{Say it:} \emph{ñān niraparādhiyāṇŭ ennŭ avan paṟaññu.} — He said that he was innocent.
  \item \emph{Recall:} say the Malayalam for the guilty one, then the Malayalam for bail, then say \emph{niraparādhi} again
\end{itemize}

\begin{tcolorbox}[breakable,colback=teal!4,colframe=teal!35!black,title={Before you move on}]
What does \ml{നിരപരാധി} mean? (\textquotedblleft{}An innocent person\textquotedblright{}.) Say it once more.
\end{tcolorbox}
\section[\texorpdfstring{\ml{വിധിക്കുക} (vidhikkuka)}{വിധിക്കുക (vidhikkuka)}]{\ml{വിധിക്കുക} (vidhikkuka) — to pass sentence; to decree}

\label{lesson:ML-C369-vidhikkuka}



\begin{quote}
Before the new one: say the Malayalam for bail, then the Malayalam for an innocent person.
\end{quote}

\subsection*{You'll want to know: \ml{വിധിക്കുക}}
\textbf{\ml{വിധിക്കുക}} — \emph{vidhikkuka} — \textquotedblleft{}to pass sentence; to decree\textquotedblright{}.

\begin{quote}
\ml{കോടതി} \ml{അവന്} \ml{വലിയ} \ml{പിഴ} \ml{വിധിച്ചു}.
\end{quote}

\begin{quote}
\emph{kōṭati avanŭ valiya piḻa vidhiccu.}
\end{quote}

\begin{quote}
The court sentenced him to a large fine.
\end{quote}

Past: \ml{വിധിച്ചു}, \emph{vidhiccu}.

The person sentenced takes -\ml{ക്ക്} or -\ml{ന്}: \ml{അവന്} \ml{പിഴ} \ml{വിധിച്ചു}.

\begin{cousinweb}
\textbf{\ml{വിധിക്കുക}} (\emph{vidhikkuka}) is built on \ml{വിധി} (\emph{vidhi}), \textquotedblleft{}a verdict\textquotedblright{}, from Sanskrit \emph{vidhi} (CDIAL 11755), with the Malayalam verb ending -\ml{ിക്കുക}: to lay down what holds in one case.
\end{cousinweb}

\begin{grammarlens}[title={What you've built}]
One more word for telling what happened.
\end{grammarlens}

\subsection*{Your turn}
\begin{itemize}
\raggedright
  \item \emph{Say it:} \emph{vidhikkuka}
  \item \emph{Say it:} \emph{vidhikkuka}, once more
  \item \emph{Say it:} \emph{kōṭati avanŭ valiya piḻa vidhiccu.} — The court sentenced him to a large fine.
  \item \emph{Recall:} say the Malayalam for bail, then the Malayalam for an innocent person, then say \emph{vidhikkuka} again
\end{itemize}

\begin{tcolorbox}[breakable,colback=teal!4,colframe=teal!35!black,title={Before you move on}]
What does \ml{വിധിക്കുക} mean? (\textquotedblleft{}To pass sentence; to decree\textquotedblright{}.) Say it once more.
\end{tcolorbox}
